\chapter{Introduction}
\label{ch:intro}

\section{The problem}
\label{sec:problem}

Early in the design of a building, someone must decide where its columns
will stand. At that point the architectural drawing exists and the
structural design does not: the drawing records walls, rooms, doors and
dimensions, but it contains no statement of which elements will carry the
weight of the floors above. The decision is usually made quickly and from
experience, and it constrains a great deal of what follows, including the
distances floors must span, the depth of the floor construction, the
number and size of foundations, and how much of the architect's intended
layout survives contact with the structure.

Two properties make this a natural target for automation. The decision is
made from information that is already digital, since architectural plans
are drawn in vector form and their wall geometry can be read directly. And
the space of reasonable answers is discrete and enumerable: a column
arrangement is a choice of grid lines together with a rule for which of
their crossings receive a column, so for a given plan and a given
permissible range of spans the candidates can be listed exhaustively or
sampled representatively.

The difficulty lies not in generating candidates but in judging them. To
say that one arrangement is better than another requires either a model of
what makes an arrangement good, which invites disagreement about the
model, or examples of arrangements that real engineers produced, which are
not present in architectural datasets. This thesis confronts that
difficulty directly rather than assuming it away, and the resulting
measurement discipline is the greater part of its contribution.

\section{Research questions}
\label{sec:questions}

The work is organised around three questions, in rising order of
importance.

\begin{enumerate}
\item \textbf{Can candidate arrangements be enumerated such that the real
building's arrangement is among them?} This must be established first. If
the correct answer is absent from the candidate set, then every accuracy
figure measures that absence rather than the quality of any method, and
the comparison that follows is meaningless. Chapter~\ref{ch:results}
therefore reports candidate coverage before any accuracy figure.

\item \textbf{Can a purely geometric score order those candidates so that
the real arrangement is near the top?} By purely geometric is meant a
score computed from the positions of the columns and the geometry of the
plan, containing no member sizes, no material prices and no load values.
The motivation for that restriction is given in
Section~\ref{sec:removals} and developed in Chapter~\ref{ch:method}.

\item \textbf{Do learned classifiers improve on that score?} Three are
tested, reading respectively numerical measurements of an arrangement, a
picture of it, and a four-qubit quantum encoding of it. All three are
trained on identical examples and labels, and the score-only pipeline is
retained as the baseline that each must beat.
\end{enumerate}

\section{Structure of the pipeline}
\label{sec:pipeline}

Figure~\ref{fig:pipeline} sets out the stages. A plan and a set of user
constraints enter at the left. The parser converts wall segments into grid
lines and candidate column positions. The generator lists the arrangements
those grid lines permit within the user's span limits. The score assigns
each candidate a single number. At that point the pipeline branches into
the four variants that the thesis compares, and each produces one chosen
arrangement. The evaluation stage measures all four against the derived
labels on plans that no model has seen.

\begin{figure}[H]
\centering
\small
\begin{tabular}{@{}rll@{}}
\multicolumn{3}{@{}l}{\textit{Input:} architectural plan (vector walls) $+$ user constraints}\\[4pt]
\midrule
Stage 04 & parser & walls $\rightarrow$ grid lines $\rightarrow$ candidate positions \\
Stage 05 & generator & every arrangement within the span limits \\
Stage 06 & score & five geometric measures $\rightarrow$ one penalty \\
 & & (weights fixed, or learned in stage 08) \\
\midrule
\multicolumn{3}{@{}l}{\textit{The branch: four variants, each choosing one arrangement}}\\[2pt]
V1 & stage 07 & re-ranked by the feature judge \\
V2 & stage 11 & re-ranked by the image judge \\
V3 & --- & not re-ranked; the score's own order stands (baseline) \\
V4 & stage 10 & re-ranked by the quantum judge \\
\midrule
Stage 12 & evaluation & coverage, accuracy, significance, sensitivity \\
\bottomrule
\end{tabular}
\caption[The pipeline and its four variants]{The pipeline and its four
variants. Stage numbers refer to the repository folders that implement each
step. All four variants share every stage before the branch, so the
comparison between them is controlled: identical grid, identical
candidates, identical score.}
\label{fig:pipeline}
\end{figure}

The three judges consume the score's fifteen best candidates and re-order
them by combining the score with their own estimate of plausibility,
\begin{equation}
\text{combined} = z(\text{penalty}) + \lambda\bigl(1 - P(\text{plausible})\bigr),
\label{eq:blend}
\end{equation}
where $z$ is a robust standardisation computed over the whole candidate
set and $\lambda$ controls how much authority the judge holds. The value
of $\lambda$ is varied in Chapter~\ref{ch:results} rather than asserted.

\section{Objectives}
\label{sec:objectives}

\begin{enumerate}
\item Derive a column-arrangement label for open floor-plan corpora,
document every assumption in the derivation, and measure how far the
results depend on the one arbitrary constant it contains.
\item Align candidate generation with that derivation until the label is
reachable, and report the coverage rate before any accuracy figure.
\item Build a score whose only free parameters are five interpretable
weights, and obtain those weights in two independent ways so that the
hand-set and learned versions can be compared.
\item Train three judges on one identical task, with implausible examples
constructed in a way that does not leak the answer.
\item Compare all four variants under a single protocol using paired
statistics, retaining the algorithm-only path as the baseline.
\item Leave every known weakness either closed or explicitly stated.
\end{enumerate}

\section{Scope, and what was deliberately removed}
\label{sec:removals}

An earlier version of this pipeline estimated a construction cost. It
selected a steel section for every beam and column, priced the steel by
weight, added a charge for each column standing away from a wall and for
each additional section size, and billed erection labour by the hour. That
machinery has been removed, and the reasoning is worth stating because it
shaped the present design.

A cost estimate is only as defensible as its unit rates. The rates in
question were indicative figures with no auditable source, and the first
question any examiner would ask of them has no good answer. What the
industry evidence behind them actually establishes is not a price but a
set of \emph{geometric drivers}: that fabrication and erection dominate
the cost of a steel package, that connections drive a large share of
fabrication cost while contributing little of the tonnage, and that
repeated bay sizes reduce cost through repeated formwork and simpler
procurement. Those drivers can be expressed as geometry and ranked without
attaching money to them, which is what the score of
Chapter~\ref{ch:method} does.

Two further removals follow from the first. Member sizing existed only to
support pricing, so it went with it; and a member-by-member check against
the Eurocodes, which verified the sections that the pricing stage had
selected, could not outlive the sizing it depended on. A lateral drift
calculation briefly replaced that check and has since also been removed.

The consequence is stated plainly rather than minimised. The pipeline
performs no structural analysis. Structural reasoning enters at exactly
one point, the span-demand measure of
Section~\ref{sec:score}, and it enters as a proportionality rather than as
a calculation. Agreement with the derived labels is therefore the only
basis on which the four variants are compared, and no measurement in this
thesis is independent of the labelling procedure. This is recorded as a
limitation in Section~\ref{sec:limitations} and as the leading item of
further work in Section~\ref{sec:further}.

\section{Contribution}
\label{sec:contribution}

The claim was checked against literature published between 2020 and 2026,
and Chapter~\ref{ch:literature} reports what was found, including work
that narrows it. No individual component of this pipeline is new.
Predicting column positions from plans with a convolutional network has
been done \citep{ampanavos2021approxiframer}, and the decomposition from
walls to grid lines to crossings is claimed in a granted patent
\citep{wang2024columnrl}. Learning structural layout from real buildings
has been done for shear walls \citep{liao2021shearwallgan,pizarro2021walls}.
Applying a variational quantum classifier to building data has been done
for post-earthquake safety assessment \citep{bhatta2024quantum}.

What was not found is any study that applies an algorithmic score,
feature-based learning, image-based learning and quantum learning to
\emph{one identical task, on one corpus, with one set of labels}, while
retaining the algorithm-only path as a baseline that the learned variants
must beat. That controlled comparison is the contribution. Two supporting
contributions are a documented and sensitivity-tested procedure for
deriving column-arrangement labels for public floor-plan datasets, and a
method for obtaining the score weights by ranking real arrangements
against alternatives drawn from the same plan.

\section{Outline of the report}

Chapter~\ref{ch:literature} reviews the relevant literature and states the
gap. Chapter~\ref{ch:method} describes the corpus, the labelling
procedure, the parser, the generator and the score.
Chapter~\ref{ch:models} describes the four models that were trained and
the design decisions that make their comparison fair.
Chapter~\ref{ch:results} reports the results, the significance tests and
the sensitivity analyses. Chapter~\ref{ch:conclusion} discusses what the
results mean, states the limitations, and identifies further work.
